\documentclass[UTF8,fontset=fandol,11pt]{ctexart}
\usepackage[a4paper,margin=25mm]{geometry}
\usepackage{amsmath,amssymb,amsthm,mathtools}
\usepackage{xcolor}
\usepackage{booktabs,array,enumitem,fancyhdr,hyperref}
\hypersetup{colorlinks=true,linkcolor=blue!50!black,urlcolor=blue!50!black}
\setlength{\emergencystretch}{2em}
\setlength{\parskip}{0.3em}
\linespread{1.16}
\allowdisplaybreaks
\newtheorem{theorem}{定理}[section]
\newtheorem{proposition}[theorem]{命题}
\newtheorem{lemma}[theorem]{引理}
\theoremstyle{definition}
\newtheorem{definition}[theorem]{定义}
\theoremstyle{remark}
\newtheorem{remark}[theorem]{说明}
\newcommand{\dd}{\,\mathrm d}
\newcommand{\ii}{\mathrm i}
\newcommand{\Rea}{\operatorname{Re}}
\newcommand{\Ima}{\operatorname{Im}}
\newcommand{\Li}{\operatorname{Li}}
\newcommand{\supp}{\operatorname{supp}}
\pagestyle{fancy}
\fancyhf{}
\fancyhead[L]{\small 学习笔记}
\fancyhead[R]{\small 初稿}
\fancyfoot[C]{\thepage}
\setlength{\headheight}{15pt}
\author{}
\date{2026年10月22日初稿}
\makeatletter
\renewcommand{\maketitle}{%
  \thispagestyle{plain}%
  \begin{center}
    {\LARGE\bfseries \@title\par}
    \vspace{0.6em}
    {\small \@date\par}
  \end{center}
  \vspace{0.6em}}
\makeatother
\newcommand{\draftnotice}{\noindent{\small
本文为供逐篇审阅的学习笔记初稿，不是新增研究成果。
文中区分直接证明与引用结果；正文不依赖本文。
本轮已作公式、条件及排版初步核查，尚待作者逐句审定。}\par\medskip}

\title{$\Gamma$ 倍元公式的两种证明}
\begin{document}
\maketitle
\draftnotice
\section{同一个公式的两条进入途径}
目标是证明
\[
 \Gamma(s)\Gamma(s+\tfrac12)=2^{1-2s}\sqrt\pi\,\Gamma(2s).
\]
先在 $\Rea s>0$ 建立等式，再说明亚纯延拓后的含义。
第一种方法从积分和变量代换进入；第二种方法从乘积的对数导数进入。
共同前提为 $\Gamma$ 的积分定义、递推关系及其在右半平面无零点的结论；
第二法还引用 Gauss 极限公式及相应的局部一致求导结果。

\section{先建立 Beta--Gamma 关系}
对 $\Rea a,\Rea b>0$，绝对收敛允许把两个积分相乘：
\[
 \Gamma(a)\Gamma(b)=
 \int_0^\infty\int_0^\infty e^{-x-y}x^{a-1}y^{b-1}\dd x\dd y.
\]
令 $x=ru$、$y=r(1-u)$，其中 $r>0$、$0<u<1$。
Jacobian 的绝对值为 $r$，故积分化为
\[
 \left(\int_0^\infty e^{-r}r^{a+b-1}\dd r\right)
 \left(\int_0^1u^{a-1}(1-u)^{b-1}\dd u\right).
\]
实正变量的幂均取实对数定义；取绝对值后同一计算仍有限，因而换元、换序对复参数合法。
于是
\[
 B(a,b)=\frac{\Gamma(a)\Gamma(b)}{\Gamma(a+b)}.
\]

\section{第一法：对称区间上的平方代换}
对 $\Rea s>0$，令 $u=(1+v)/2$，得到
\begin{align*}
 B(s,s)
 &=2^{1-2s}\int_{-1}^1(1-v^2)^{s-1}\dd v\\
 &=2^{2-2s}\int_0^1(1-v^2)^{s-1}\dd v\\
 &=2^{1-2s}\int_0^1w^{-1/2}(1-w)^{s-1}\dd w\\
 &=2^{1-2s}B(\tfrac12,s).
\end{align*}
第三行用了 $w=v^2$、$\dd v=\tfrac12w^{-1/2}\dd w$。
将 Beta--Gamma 关系代入，并约去无零点的 $\Gamma(s)$，便得
\[
 \frac{\Gamma(s)}{\Gamma(2s)}
 =2^{1-2s}\frac{\Gamma(1/2)}{\Gamma(s+1/2)}.
\]
这里 $\Gamma(1/2)=\sqrt\pi$：令 $x=t^2$ 可化为 Gaussian 积分，
而 $\int_{\mathbb R}e^{-t^2}\dd t=\sqrt\pi$ 可由其平方在平面上作极坐标积分得到。
具体地，
$\int_{\mathbb R^2}e^{-x^2-y^2}\dd x\dd y
=2\pi\int_0^\infty re^{-r^2}\dd r=\pi$，正平方根确定符号。
第一法完成。

\section{第二法：把整数与半整数两列配在一起}
为避免与 Chebyshev 函数 $\psi$ 混淆，记
$D(s)=\Gamma'(s)/\Gamma(s)$。
由 Gauss 极限公式及其在紧子集上的对数导数收敛，
\[
 D(s)=\lim_{N\to\infty}
 \left(\log N-\sum_{k=0}^{N}\frac1{s+k}\right),
 \qquad\Rea s>0.
\]
本节引用这一局部一致极限；它可由 Gauss 乘积的通常证明获得，不以倍元公式为前提。

有限和有精确关系
\[
 \sum_{k=0}^{N}\left(\frac1{s+k}+\frac1{s+1/2+k}\right)
 =2\sum_{j=0}^{2N+1}\frac1{2s+j}.
\]
对右侧的 $D(2s)$ 取截断指标 $2N+1$，因而
\begin{align*}
 D(s)+D(s+\tfrac12)-2D(2s)
 &=\lim_{N\to\infty}\bigl(2\log N-2\log(2N+1)\bigr)\\
 &=-2\log2.
\end{align*}
现在令
\[
 Q(s)=2^{2s-1}\frac{\Gamma(s)\Gamma(s+1/2)}{\Gamma(2s)}.
\]
它在右半平面解析且无零点，并满足
$Q'(s)/Q(s)=2\log2+D(s)+D(s+1/2)-2D(2s)=0$。
所以 $Q$ 在连通的右半平面上为常数。
取 $s=1/2$，用 $\Gamma(1)=1$ 得 $Q(1/2)=\sqrt\pi$，再次得到目标公式。

\section{延拓、极点与方法比较}
两边都是全平面上的亚纯函数，并已在开集 $\Rea s>0$ 相等，故由唯一性作为亚纯恒等式成立。
若某个参数使因子有极点，应比较 Laurent 展开或极限，不把“无穷大乘无穷大”当作普通数值代入。

例如 $s=1$ 时，
$\Gamma(3/2)=\sqrt\pi/2$，与递推关系一致。
第一法显示对称与平方代换如何产生半整数；第二法显示整数步长的两列如何合并成半步长的一列。
二者证明同一公式，但组织信息的方式不同。
第一法所需的是绝对收敛积分的换元与换序；第二法则把主要分析前提集中在
Gauss 极限的局部一致求导上。后者的代数计算较短，并不意味着其分析前提更少。

\section*{参考依据与后续核验}
Beta 积分及倍元公式的标准条目可查 NIST DLMF，§5.12 与 §5.5：
\url{https://dlmf.nist.gov/5.12}，
\url{https://dlmf.nist.gov/5.5}。
本书《Gamma 函数及其定义域的扩展》提供 Gauss 极限及相关延拓背景。
本文已明确第二法引用的局部一致求导结果，后续审阅应特别核对这一依赖是否需要另附证明。
\end{document}
